\chapter{天气接口：数据库优先、外部兜底，以及时间匹配}

最后把天气提供给前端，并且准备好“预测用真实天气”所需的匹配逻辑。

\partinfo{%
  \code{weather/WeatherQueryService.java}：13.1\\
  \code{weather/web/WeatherController.java}、\code{WeatherExceptionHandler.java}：13.2\\
  \code{weather/WeatherMatcher.java}：13.3%
}{\code{./mvnw test -Dtest='Part13*'}，最后 \code{./mvnw test}}

\section{知识点}

\subsection{数据库优先，外部 API 兜底}

\begin{plaincode}
查天气 --> 数据库里有？ --有--> 用数据库的（快，不消耗 API 调用次数）
                  |
                  没有（刚启动、抓取失败过）
                  v
            调用 OpenWeather（有 5 分钟缓存）--> 也失败 --> WeatherUnavailableException --> 503
\end{plaincode}

原项目前端自己做了这套逻辑：先请求一个接口，失败了再请求另一个。现在挪到后端，前端只要请求一个接口，以后要改兜底的顺序，也不用改前端。

\begin{note}[同一个 \code{time()}，两种意思]
数据库的两个实体都有 \code{toSnapshot()}，OpenWeather 客户端返回的也是 \code{WeatherSnapshot}。但 \code{time()} 的含义要看它是哪种数据：
\begin{itemize}
  \item \textbf{实时天气}的 \code{time()} = \code{observedAt}：OpenWeather \textbf{测量}的时间，一般是几分钟前（Part 12）。
  \item \textbf{预报}的 \code{time()} = \code{forecastTime}：这条预报说的是\textbf{未来哪个时刻}。
\end{itemize}
13.1b 里的“现在”是一条\textbf{分界线}：只留 \code{time()} 不早于现在的预报，已经过去的时刻扔掉。它不是用来找“等于现在”的那条。
\end{note}

\subsection{多个 @RestControllerAdvice}

Part 6 写过 \code{GlobalExceptionHandler}，它继承了 \code{ResponseEntityExceptionHandler}，负责 Spring 自带的那些异常。天气的异常可以单独写一个 advice 放在天气模块里，模块之间互不干扰。

\textbf{但只能有一个 advice 继承 \code{ResponseEntityExceptionHandler}}：它里面已经给十几种 Spring 自带的异常（参数校验失败、请求体读不了、缺少参数……）写好了 \code{@ExceptionHandler}。两个类都继承的话，这些异常就有了两个处理器。Spring \textbf{不会报错}，而是按 advice 的顺序悄悄选一个，选中哪个取决于 Bean 的排序，很难预料。后果是：你在天气模块里改了校验失败的响应，车站接口的 400 也可能跟着变（我故意试过：在 \code{WeatherExceptionHandler} 里把校验失败改成 418，结果 Part 6 车站接口的校验测试也收到了 418）。所以这类“Spring 自带异常”的处理只放在一个地方（Part 6 的 \code{GlobalExceptionHandler}），各模块的 advice 只处理自己模块的异常。

\subsection{本地时间和 Instant}

预测用的时间点是 \code{LocalDateTime}（都柏林的“几点”，Part 9），天气预报的时间是 \code{Instant}（UTC 的一个时刻）。比较前要先换算：

\begin{java}
Instant t = localDateTime.atZone(ZoneId.of("Europe/Dublin")).toInstant();
Duration gap = Duration.between(a, b).abs();          // 两个时刻相差多久（取绝对值）
\end{java}

冬天都柏林和 UTC 一样，夏天快 1 小时。测试里有一个 7 月的例子专门检查这一点：如果直接用 \code{ZoneOffset.UTC} 换算，冬天的测试能过，夏天的会挂。

13.3 里要对上的两种时间：

\begin{center}
\begin{tabularx}{\textwidth}{l l >{\raggedright\arraybackslash}X}
\toprule
参数 & 类型 & 是什么 \\
\midrule
\code{times} & \code{LocalDateTime} & 要预测的那些整点，\textbf{都柏林本地时间}，就是 Part 9 \code{forecastTimes(hours)} 算出来的 \\
\code{forecast} 里每条的 \code{time()} & \code{Instant} & 这一档预报说的时刻（UTC），每 3 小时一档 \\
\code{zone} & \code{ZoneId} & 用来把 \code{times} 换算成 \code{Instant}；调用方传 \code{clock.getZone()}，测试可以传别的时区 \\
\bottomrule
\end{tabularx}
\end{center}

“最接近”指两个 \code{Instant} 相差的时间（取绝对值）最小。预报是 3 小时一档，整点是 1 小时一个，所以几个连续的整点会对到同一档预报，这是正常的。

“找最接近的那一个”可以用 Stream 的 \code{min} 加 \code{Comparator}：

\begin{java}
Optional<Shop> closest = shops.stream()
        .min(Comparator.comparing((Shop s) -> distanceTo(s))     // 先按距离
                .thenComparing(Shop::name));                     // 距离一样时按名字
\end{java}

\section{题目}

\begin{task}{13.1\quad WeatherQueryService}
Service，所有方法在只读事务里。

\textbf{a) \code{current()}}：数据库里最新的一次观测（Part 12 的 \code{latest()}，再用实体的 \code{toSnapshot()} 转成 \code{WeatherSnapshot}）；数据库是空的（\code{Optional} 为空）就调用 \code{client.current()}。

\textbf{b) \code{forecast(limit)}}：先查数据库 \code{upcoming(现在, limit)}，“现在”用 \code{Instant.now(clock)}，每条用 \code{toSnapshot()} 转换；一条都没有，就用 \code{client.forecast()}。OpenWeather 返回的列表\textbf{可能包含已经过去的那一档}（下面例子里的 15:00），所以同样要过滤：只留 \code{time()} 不早于“现在”的，再取最多 \code{limit} 条。

\begin{example}
\begin{plaincode}
现在 18:05，数据库没有预报，OpenWeather 返回 15:00, 21:00, 00:00, 03:00
forecast(2)  ->  [21:00, 00:00]
\end{plaincode}
\end{example}
\end{task}

\begin{task}{13.2\quad WeatherController 和 503}
\textbf{a)} REST Controller，公共路径 \code{/api/v1/weather}：
\begin{itemize}
  \item \code{GET /current} → \code{WeatherResponse.from(...)}
  \item \code{GET /forecast?limit=..} → 列表，每条也用 \code{WeatherResponse.from}；\code{limit} 默认 8，范围 1--40。写法和 Part 6 的 \code{limit} 一样：\code{@RequestParam(defaultValue = ...)} 加 \code{@Min} / \code{@Max}。超出范围时 Spring 抛校验异常，Part 6 的 \code{GlobalExceptionHandler} 会把它转成 400，这里不用再处理。
\end{itemize}

\textbf{b)} \code{WeatherExceptionHandler}：\code{WeatherUnavailableException} 转成 503 ProblemDetail，title 是 \code{"Weather unavailable"}，detail 是异常消息。写法参照 Part 6 \code{GlobalExceptionHandler} 里处理 404 的那个方法，只是状态码换成 \code{HttpStatus.SERVICE_UNAVAILABLE}，并且这个类\textbf{不要}继承 \code{ResponseEntityExceptionHandler}（原因见知识点）。
\end{task}

\begin{task}{13.3\quad WeatherMatcher（先看测试再写）}
\code{match(times, forecast, zone)}：对 \code{times} 里的每个本地时间点，先用 \code{zone} 换算成 \code{Instant}，再在 \code{forecast} 里找 \code{time()} 最接近的那一档，转成 \code{WeatherInput}（Part 7 的 record）。返回的列表和 \code{times} 一样长、顺序一一对应：
\begin{itemize}
  \item 温度、湿度、气压照搬；
  \item 下雨：\code{rainMm > 0.01} 或者 \code{rainProbability >= 0.5} 时 \code{isRaining = 1}，否则 0；
  \item 前后两档一样近时，取\textbf{更早}的那一档；
  \item \code{forecast} 为空：返回 \code{null}（Python 服务收到没有 weather 的请求，会自己推断天气）。
\end{itemize}

\begin{example}
\begin{plaincode}
预报：18:00（9 度）、21:00（8 度，降雨 0.9mm）、00:00（7 度）
时间点 18:00 -> 18:00 那档，不下雨
时间点 20:00 -> 21:00 那档，下雨
时间点 19:30 -> 离 18:00 和 21:00 一样近，取 18:00
\end{plaincode}
\end{example}
\end{task}

\section{可选：让预测用上真实天气}

\code{WeatherMatcher} 准备好了，但 Part 9 的 \code{StationForecastService} 还在传 \code{weather = null}。想接上的话：给它的构造函数加一个 \code{WeatherQueryService} 参数，在 \code{forecast()} 里

\begin{java}
List<WeatherInput> weather = WeatherMatcher.match(times, weatherQueryService.forecast(16), clock.getZone());
ForecastRequest request = new ForecastRequest(number, station.totalStands(), times, weather);
\end{java}

为什么是 16：预测最多 48 小时（Part 9 的 \code{hours} 上限），预报 3 小时一档，$48 \div 3 = 16$ 档正好覆盖。

这样会改变 Part 9 测试里用到的构造函数，所以要同时把 \code{Part9ForecastServiceTest} 里的 \code{new StationForecastService(...)} 补上一个参数（可以传 \code{mock(WeatherQueryService.class)}）。这一步没有专门的测试，做完用真实数据跑一下，比较接上前后的预测结果。

\begin{think}
\begin{enumerate}
  \item “数据库优先、API 兜底”和“API 优先、数据库兜底”，各适合什么场景？
  \item 为什么 \code{LocalDateTime} 和 \code{Instant} 不能直接比较？夏令时切换的那一天，本地时间 01:30 可能出现两次，会有什么影响？
  \item 天气接口失败返回 503，而预测接口失败返回 200 + \code{degraded}，为什么不一样？
\end{enumerate}
\end{think}
